%% kapa: exam=1 round=27 year=2016 period=1 form=A subject=real_estate subject_ko=부동산학원론 seq=25 number=105
%% source: Q-Net 감정평가사 기출문제 (cst00309) artlSeq=5169402 file=extracted/1008/5169402/제27회 감정평가사 1차 문제/제27회 감정평가사 1교시문제(A형).pdf
%% conversion: kordoc 4.24.51 -> tools/gen1.py (auto); review=proofread-sonnet; flags=none
\begin{question}{105}
감정평가에 관한 규칙상 감정평가에 관한 설명으로 옳지 않은 것은?
\begin{choices}
  \item 토지를 감정평가할 때에 부동산 가격공시 및 감정평가에 관한 법률에 따라 공시지가기준법을 적용하여야 한다.
  \item 공시지가기준법에 따라 토지를 감정평가할 때에는 비교표준지 선정, 시점수정, 지역요인 비교, 개별요인 비교, 그 밖의 요인 보정의 순서에 따라야 한다.
  \item 건물을 감정평가할 때에 원가법을 원칙적으로 적용하여야 한다.
  \item 과수원을 감정평가할 때에 수익환원법을 원칙적으로 적용하여야 한다.
  \item 자동차를 감정평가할 때에 거래사례비교법을 원칙적으로 적용하여야 하나, 본래 용도의 효용가치가 없는 물건은 해체처분가액으로 감정평가할 수 있다.
\end{choices}
\end{question}
\answer{4}
